\chapter{The financial model and credit mapping}
\label{ch:finance}
\modtag{M5/M9}{Financial model and credit}{Deterministic core}

This is the chapter where a battery and its dispatch schedule turn into something a lender can read. Everything upstream answered ``what should we build and how should it run''; the financial model answers ``does it pay, who lends against it, and at what price''. It takes the chosen system, the value its dispatch earns, the scenario cases, and the customer's credit standing, and it produces the numbers that fill the underwriting pack: cashflows in a normal year and in a bad year, whether those cashflows cover the loan, how much of the income is contracted rather than speculative, and the price the project must charge to hit its target return.

Two design promises govern every line of it. First, \emph{no artificial intelligence touches any number here}: this module is pure arithmetic, the same inputs always give the same answer, and every figure is read either from the validated site or from the one versioned assumption library. Second, the module \emph{degrades honestly}. When an input is missing it produces an explicitly \emph{unavailable} field and a warning, never a guessed stand-in; and when there is no positive system to price at all, it refuses outright rather than inventing one.

\begin{plainnote}[A quick finance glossary]
A few terms recur throughout. \textbf{Cashflow}: the cash a project earns in a year after running costs, before any loan repayment. \textbf{P50 / P90}: two views of that cashflow -- P50 is the central, ``expected'' case; P90 is a deliberately cautious case (the level you would beat nine years in ten). Lenders underwrite to P90. \textbf{Capex}: the up-front cost of buying and installing the battery. \textbf{Debt service}: the annual loan repayment (interest plus principal). \textbf{DSCR} (debt service cover ratio): cashflow divided by debt service in a year; 1.0 means the cash exactly covers the repayment, higher is safer. \textbf{Gearing}: the share of capex funded by a loan rather than equity. \textbf{Tenor}: the length of the loan in years. \textbf{Sculpting}: shaping the repayment so the cover ratio stays at a chosen level rather than letting it drift. \textbf{IRR}: the annualised return a cashflow stream earns. These are all defined again in Appendix~\ref{app:glossary}.
\end{plainnote}

\pipelinepos{receives the chosen system + valued dispatch from the sizing optimiser (Ch.~\ref{ch:sizing}) and dispatch optimiser (Ch.~\ref{ch:dispatch}), the Low/Central/High cases from the scenario engine (Ch.~\ref{ch:scenarios}), and the versioned numbers from the assumption library (Ch.~\ref{ch:assumptions}); hands a \code{FinancialCase} to the quote assembler (Ch.~\ref{ch:quotes}).}

\begin{contract}
\textbf{In:} a validated \code{SiteModel} (its tariff levy rates are read directly); an \code{OptimisationResult} or any compatible result object carrying the chosen power, energy and duration plus the per-stream value and firmness; the versioned \code{AssumptionLibrary} (the only source of every cost, rate and ratio); optionally a scenario (price and cost multipliers), an annual dispatch energy accounting (charge / discharge kWh), an injected debt structure, and a choice of commercial structure.

\smallskip
\textbf{Out:} a \code{FinancialCase} record carrying the P50 and P90 cashflow series, the DSCR summary (minimum and by-year, measured on P90), project IRR and levered equity IRR, simple payback, the debt assumptions actually used (gearing, rate, tenor, sculpting basis), the price recommendation, the covenant verdict, and a deduplicated list of warnings. Alongside it the module returns the full cashflow detail, the DSCR detail and a \code{RunContext} provenance stamp, and a single \code{bankable} flag.

\smallskip
\textbf{Refuses:} when there is no positive system to price (no chosen power), it raises an ``insufficient inputs'' error so the caller produces a ranked data ask, not a priced fiction.
\end{contract}

\section{The shape of the calculation}

The model runs as a fixed, ordered pipeline. Each stage consumes the one before it, and nothing reaches back upstream, so the whole thing is one deterministic pass:

\begin{center}\sffamily\small
cashflows $\;\rightarrow\;$ P50/P90 series $\;\rightarrow\;$ debt and DSCR $\;\rightarrow\;$ payback and IRR $\;\rightarrow\;$ price $\;\rightarrow\;$ \code{FinancialCase}
\end{center}

Figure~\ref{fig:finance-flow} is the full map. Read it top to bottom: the four inputs on the left feed the cashflow builder, the cashflows split into the P50 and P90 views, debt is sized and measured on the cautious P90 view, returns are read off the central P50 view, a price is solved, and everything is assembled and stamped into the record.

\begin{figure}[H]
\centering
\resizebox{\textwidth}{!}{%
\begin{tikzpicture}[node distance=7mm and 15mm]
  % ---- inputs on the left ----
  \node[io] (optr) {\textbf{OptimisationResult}\\\scriptsize chosen size + per-stream\\\scriptsize value \& firmness};
  \node[io, below=of optr] (scen) {\textbf{Scenario}\\\scriptsize price \& cost multipliers\\\scriptsize (Low/Central/High)};
  \node[io, below=of scen] (assum) {\textbf{AssumptionLibrary}\\\scriptsize every \pounds/\%/ratio,\\\scriptsize versioned};
  \node[io, below=of assum] (credit) {\textbf{Credit status}\\\scriptsize from \code{SiteModel}\\\scriptsize counterparty band};

  % ---- spine ----
  \node[proc, right=24mm of scen, minimum width=40mm, text width=38mm] (cf)
     {\textbf{Build cashflows}\\\scriptsize revenue by firm stream,\\\scriptsize cost lines, capex};
  \node[decide, left=14mm of cf] (gate) {sized?};

  \node[proc, right=of cf, minimum width=38mm, text width=36mm] (pxx)
     {\textbf{P50 / P90 split}\\\scriptsize haircut + degradation\\\scriptsize decay per year};

  \node[proc, below=14mm of pxx, minimum width=38mm, text width=36mm] (debt)
     {\textbf{Size \& measure debt}\\\scriptsize on the P90 case\\\scriptsize $\rightarrow$ DSCR series + min};
  \node[ext, left=of debt, minimum width=34mm, text width=32mm] (map)
     {\textbf{Credit $\rightarrow$ rate/tenor}\\\scriptsize M9 deterministic map};

  \node[proc, below=10mm of cf, minimum width=38mm, text width=36mm] (ret)
     {\textbf{Returns}\\\scriptsize payback, project IRR\\\scriptsize (P50); equity IRR};
  \node[proc, below=10mm of ret, minimum width=38mm, text width=36mm] (split)
     {\textbf{Contracted vs merchant}\\\scriptsize firm/contracted floor\\\scriptsize vs conditional upside};
  \node[proc, right=of ret, minimum width=38mm, text width=36mm] (price)
     {\textbf{Price-to-target}\\\scriptsize \pounds/kWh that hits the\\\scriptsize required return};

  \node[store, below=12mm of price, minimum width=42mm, text width=40mm] (case)
     {\textbf{FinancialCase}\\\scriptsize cashflows, DSCR, IRR,\\\scriptsize price, covenant, warnings};
  \node[term, right=14mm of case] (out) {to quote\\record};

  % arrows in
  \draw[flow] (optr.east) to[out=0,in=180] (gate.west);
  \draw[flow] (scen.east) to[out=0,in=160] (gate.north west);
  \draw[flow] (assum.east) to[out=0,in=200] (gate.south west);
  \draw[flow] (gate) -- node[flowlabel,above]{yes} (cf);
  \draw[flow] (credit.east) to[out=0,in=180] (map.west);

  % spine
  \draw[flow] (cf) -- (pxx);
  \draw[flow] (pxx) -- (debt);
  \draw[flow] (map) -- (debt);
  \draw[flow] (cf) -- (ret);
  \draw[flow] (cf.south) to[out=-90,in=90] (split.north);
  \draw[flow] (ret) -- (price);
  \draw[flow] (pxx.south) to[out=-90,in=90] (price.north);
  \draw[flow] (debt) to[out=0,in=180] (price.west);
  \draw[flow] (price) -- (case);
  \draw[flow] (ret.south) to[out=-90,in=180] (case.west);
  \draw[flow] (split.east) to[out=0,in=200] (case.south west);
  \draw[flow] (case) -- (out);

  % gate fail
  \node[term, below=8mm of gate, fill=nodeIO, draw=nodeIOBd] (ask) {ranked\\data ask};
  \draw[flow] (gate.south) to[out=-90,in=90] node[flowlabel,right]{no} (ask.north);

  % zone backdrop
  \begin{scope}[on background layer]
    \node[groupbox, fit=(cf)(pxx)(debt)(ret)(split)(price)(case), inner sep=8pt,
          label={[note,rotate=90,anchor=south]west:DETERMINISTIC, NO AI}] {};
    \node[groupbox, draw=nodeExtBd, fill=gridqVioletBg, fit=(map), inner sep=4pt] {};
  \end{scope}
\end{tikzpicture}}
\caption{The financial model end to end. The whole spine is pure arithmetic; the only perimeter-flavoured box is the credit-to-rate mapping, and even that is a deterministic table lookup (the violet outline marks where the credit \emph{band} originated, not any AI in the maths). Debt is sized on the cautious P90 case; returns are read off the central P50 case.}
\label{fig:finance-flow}
\end{figure}

\section{Building the cashflows}

The first stage turns the chosen system into a year-by-year cash picture. Three things are assembled: how much the battery earns, what it costs to install, and what it costs to run.

\subsection{The gate: is there a system to price?}

Before anything else the model checks that the optimisation result actually chose a positive battery power. If it did not (an unsized or infeasible result), the model raises an ``insufficient inputs'' error. This is the honest-degradation rule in its starkest form: with nothing to size, the correct output is a \emph{ranked list of the data that would let us answer}, produced by the caller, not a fabricated system. The finance layer never invents a battery to have something to price.

\subsection{Revenue: only the firm streams count}

Revenue is read straight from the per-stream values the optimiser produced; the finance layer invents nothing, so a zero stays a zero. But not every stream is counted. Each value stream carries a \emph{firmness} label, and only the firm part forms the priced floor. The rule is: a stream counts toward priced revenue if and only if its firmness is \emph{firm} or \emph{contracted}. Conditional streams, and the separate ``gated upside'' streams the optimiser keeps in their own list (wholesale, balancing, ancillary services), are held out of the floor entirely. They are genuine potential income, but they are speculative, so the underwritten case never leans on them. Counting them would overstate both the revenue and its associated dispatch cost the moment a real engine starts pricing conditional streams.

If there is no priced revenue at all (every stream is zero or excluded, as is the case with the integration-only baseline engine), the model still proceeds, but it emits a warning and produces a \emph{cost-only}, not-bankable view. That is exactly the situation today: a real dispatch optimiser will fill in the revenue, and everything downstream follows with no code change.

\subsection{Capex from a versioned placeholder}

The installed capital cost is the chosen power, in megawatts, multiplied by an installed-cost rate in pounds per megawatt read from the library, and then by the scenario's cost multiplier. There is no hardcoded capex anywhere; the rate is one named, versioned assumption. The model rejects a negative installed-cost assumption, a negative power and a negative resulting capex, because capex can never be negative.

\begin{caution}[The capex sanity check never changes the number]
The brief requires installer capex to be checked against a benchmark. The model does exactly that: if the unit cost falls outside the library's plausibility band (a low / high in pounds per megawatt) it raises a warning telling the underwriter to review the quote. But it never alters the figure and never blocks the run. A benchmark band is a guard rail, not a price input. If no band is configured, the check is skipped rather than a band being guessed.
\end{caution}

\subsection{The cost lines}

Operating costs are built one line at a time, each from a single named assumption (or, for the levies, from the site's own tariff). The cost-line set is fixed and ordered so the result is deterministic. The crucial discipline is that a line whose driving input is missing is set to an explicit \emph{unavailable} value with a warning, never silently to zero and never invented.

\begin{table}[H]
\centering
\caption{The operating cost lines, each traceable to one source. ``Unavailable'' lines are omitted with a warning, never estimated.}
\label{tab:costlines}
\small\sffamily
\begin{tabularx}{\textwidth}{@{}l L L@{}}
\toprule
\textbf{Cost line} & \textbf{Source} & \textbf{How it is computed} \\
\midrule
Fixed opex & \fpath{battery\_opex.\seqsplit{annual\_fixed\_opex}} (\pounds/MW/yr) & rate $\times$ power (MW) \\
Insurance & \fpath{insurance.\seqsplit{annual\_pct\_of\_capex}} (\%) & percentage $\times$ capex \\
Metering / IDNO & \fpath{metering\_idno.\seqsplit{annual\_gbp\_per\_site}} & flat per site \\
DNO costs & \fpath{dno\_costs.\seqsplit{annual\_gbp\_per\_site}} & flat per site \\
Business rates & \fpath{business\_rates.\seqsplit{annual\_gbp\_per\_mw}} (\pounds/MW/yr) & rate $\times$ power (MW) \\
Levies on charge & site \fpath{tariff\_stack.levies} (\pounds/kWh) & levy rate $\times$ annual charge kWh \\
Degradation / warranty proxy & \fpath{degradation.\seqsplit{annual\_capacity\_fade\_pct}} (\%) & fade \% $\times$ capex \\
Degradation marginal cost & \fpath{battery\_opex.\seqsplit{degradation\_marginal\_cost\_gbp\_per\_mwh}} (\pounds/MWh) & rate $\times$ annual discharge MWh (release-gated) \\
\bottomrule
\end{tabularx}
\end{table}

Three lines deserve a closer look. The \textbf{levies on charge volume} line is the cost of the taxes and policy charges that apply to every unit of electricity the battery imports to charge, \emph{including the energy lost to round-trip inefficiency} (you pay the levy on what you draw from the grid, not on what you later discharge). The levy rate is the sum of the site's own tariff levy lines, which are site data, not an assumption. The annual charge volume, though, is a dispatch output that the current result schema does not yet carry, so this line stays unavailable until a dispatch energy accounting is supplied. When only the discharge volume is known, it is grossed up to a charge volume by dividing by the round-trip-efficiency assumption (and the model warns that it did so). The \textbf{degradation / warranty proxy} is a deliberate simplification: a single annual percentage of capex held as a replacement reserve, standing in for the proper ex-post battery-health and warranty model that lives in the calibration module. It is labelled as a proxy wherever it appears. The \textbf{degradation marginal cost} line prices battery wear per unit of throughput: a \pounds/MWh rate from the library multiplied by the annual discharge volume. It is \emph{release-gated}: the early releases carry no such rate (the cycling cost is carried entirely by the warranty throughput cap there, and the line stays unavailable, never zero), while from \code{gb\_v1\_1} the key exists as a labelled placeholder and the line is priced whenever a dispatch discharge volume is known (when only the charge volume is known, discharge is derived as charge $\times$ round-trip efficiency, with a warning). A negative wear rate is rejected outright, because a wear cost that pays the project would flatter the case.

\section{From cashflows to P50 and P90}

A lender does not underwrite the expected case; they underwrite a cautious one. So the model produces two parallel cashflow series from the same building blocks.

The transformation is a \emph{deterministic, assumption-driven haircut}, not a random simulation. There is no Monte Carlo here -- stochastic risk modelling belongs to the portfolio workstream, behind the same interface. For each operating year:

\begin{enumerate}
  \item the central (\textbf{P50}) revenue is the basis-year priced revenue, scaled by the scenario's price multiplier, then decayed by the capacity-fade percentage compounded over the years (a simple proxy for the battery delivering a little less each year);
  \item the cautious (\textbf{P90}) revenue is the P50 revenue multiplied by one minus a revenue-haircut percentage read from the library;
  \item costs are identical in both views (the haircut is a confidence haircut on \emph{revenue}, not on costs), so the EBITDA-like operating cashflow in each year is that year's revenue less the dispatch-level cost less the operating cost lines.
\end{enumerate}

Each series is then prefixed with year zero, which is the negative capex, to give the free cashflow before debt. The model applies no tax and no working-capital movement in this version (both documented simplifications), and applies no inflation indexation because the library carries no indexation assumption and the model never invents one. By construction the P90 series is always at or below the P50 series, and the two coincide only when revenue is zero. The model enforces that inequality structurally: a negative haircut assumption (which would lift the cautious case above the central one) is clamped to zero with a loud warning, never applied silently.

\begin{keyidea}[P90 is the underwriting case]
The split is the heart of bankability. Debt is sized and the cover ratio is measured on the \emph{P90} (cautious) cashflows, so the loan must survive a bad-but-plausible year. The headline equity return is reported on the \emph{P50} (central) cashflows, the expected outcome. This ``size to the downside, report the upside'' split is the standard discipline of project finance, and the model enforces it structurally.
\end{keyidea}

\section{Sizing the debt and the DSCR series}

With a P90 cashflow profile in hand, the model can ask the lender's question: is there cash to cover a loan, and by how much?

A debt structure has three numbers: \emph{gearing} (what fraction of capex is borrowed), the \emph{interest rate}, and the \emph{tenor} (the number of years over which it is repaid). The rate and tenor come from the credit mapping described in the next section. The gearing is portfolio-finance content: the early releases (\code{gb\_v0\_1}, \code{gb\_v1\_0}) carry none, so on them the model reports the DSCR as \emph{unavailable}, with a warning naming exactly what is missing -- it never guesses a gearing to manufacture a cover ratio. From \code{gb\_v1\_1} the library carries a gearing as a clearly labelled placeholder (illustrative, pending the portfolio-finance workstream), so the full debt structure assembles, the DSCR computes and the covenant verdict is genuinely produced end to end -- while any case priced on that placeholder stays firmly not bankable.

When a full debt structure \emph{is} present (injected by the caller or assembled from the library), the maths runs as follows:

\begin{enumerate}
  \item the loan \textbf{principal} is the capex times the gearing percentage;
  \item the annual \textbf{debt service} takes the shape named by the structure's \emph{sculpting basis}: an annuity by default (the same payment every year for the tenor, computed by the standard closed-form annuity formula, with the zero-interest case handled as a straight-line repayment), or a schedule sculpted to the P90 cashflow profile (described in the note below);
  \item the \textbf{DSCR for each year} is that year's P90 cashflow divided by that year's debt service; a year with no debt service is skipped so there is never a division by zero;
  \item the \textbf{minimum DSCR} across the debt years is the headline figure, the one the covenant tests.
\end{enumerate}

If the projection horizon is shorter than the loan tenor, the model measures the DSCR only over the overlapping years and warns that it did so, rather than padding the cashflow with invented years.

\begin{plainnote}[Sculpting, and how this version does it]
``Sculpting'' a loan means shaping the repayments year by year so the cover ratio stays pinned at a chosen level rather than drifting up and down as cashflows change; it lets a project borrow as much as the cashflows safely allow. The model supports two debt-service shapes, selected by the debt structure's \code{sculpting\_basis}. The default, \textbf{annuity}, is the level-payment schedule described above: the cover ratio is \emph{measured} and moves with the cashflow profile. The \textbf{P90-DSCR sculpt} shapes the service to the cautious cashflows directly: a flat target cover ratio $x$ is solved in closed form so that the present value, at the debt rate, of each year's P90 cashflow divided by $x$ over the tenor equals the principal; each year's service is then that year's P90 cashflow divided by $x$. By construction the discounted service repays exactly the principal and the cover ratio sits at $x$ in every debt year. The solve is deterministic and handles the zero-rate case with no special treatment. It degrades honestly: when any P90 cashflow year inside the tenor is zero or negative, or the horizon is shorter than the tenor, no positive sculpted service exists, so the model falls back to the annuity with a warning -- the service is never negative and never guessed. The record carries the basis that was \emph{actually applied}, so a requested sculpt that degraded to the annuity is recorded as the annuity, never as a sculpt that did not happen.
\end{plainnote}

The covenant itself is the minimum cover ratio a lender will accept. That value \emph{does} live in the library (unlike the gearing/rate/tenor), so the model reads it and produces a verdict: the covenant \emph{passes} when the minimum DSCR is at least the covenant. If either the measured DSCR or the covenant value is missing, the verdict is recorded as \emph{not evaluated} (neither pass nor fail), again with a warning, rather than a fabricated pass.

\section{Returns: payback and IRR}

Two return metrics are read off the central P50 cashflows, and both fail safe rather than mislead.

\textbf{Simple payback} is the year at which the running cumulative cashflow first turns positive, interpolated within the crossing year. If the project never recovers its outlay within the horizon, payback is reported as ``does not pay back'', not as a large but technically-positive number.

\textbf{Project IRR} is the discount rate at which the cashflow's net present value is zero, found by bisection. Crucially, it is computed \emph{only when a single, trustworthy answer exists}: the cashflow must have exactly one sign change (the conventional ``spend up front, earn thereafter'' shape), which guarantees a unique root. A cashflow with several sign changes can have several IRRs, none of them meaningful on its own, so the model declines to report one and points the reader to NPV and DSCR instead.

The model also reports a \textbf{levered (equity) IRR}: the return to the equity investor specifically, who funds only the capex not covered by the loan and receives each year's cashflow after debt service. It is computed on the P50 case, net of the very same debt service the P90 DSCR was measured against (the same numbers, never recomputed, so the two can never silently drift), and it shares the single-sign-change reliability gate. It is reported as unavailable when there is no debt (so there is no leverage to measure), the equity outlay is non-positive, or the series is non-conventional.

\section{Commercial structures as plug-ins}

How the customer pays is a separate choice from how the project earns. The model supports three commercial structures behind one interface, so a new structure is a new plug-in rather than a rewrite of the engine.

\begin{table}[H]
\centering
\caption{The three commercial structures. Only the fixed price is computed in this version; the other two are explicit, labelled placeholders.}
\label{tab:structures}
\small\sffamily
\begin{tabularx}{\textwidth}{@{}l L L@{}}
\toprule
\textbf{Structure} & \textbf{What the customer pays} & \textbf{Status} \\
\midrule
Fixed \pounds/kWh (default) & a fixed price per unit of energy the battery delivers them & computed (when delivered energy is known) \\
Savings share & a share of the bill savings the battery creates & labelled placeholder; needs a measurement-and-verification method first \\
Capacity fee & a fee per kW of battery capacity made available & labelled placeholder; needs a sourced \pounds/kW/yr basis \\
\bottomrule
\end{tabularx}
\end{table}

Each structure carries the commercial terms a contract needs: the \emph{term} (how many years the deal runs), \emph{indexation} (whether the price rises with inflation each year), and a \emph{buyout} (the option for the customer to buy the asset out). These terms live with the contract model (Chapter~\ref{ch:contracts}); the finance layer's job is to set the headline price. The savings-share and capacity-fee structures return their shape but no value, each with a warning saying plainly why no number is computed. This is the no-unlabelled-stand-in rule: a structure that cannot yet be priced says so, rather than emitting a placeholder figure that might be mistaken for a real one.

\subsection{The price-to-target-return solve}

For the default fixed-price structure, the model solves for the price that makes the project hit its required return. The logic is a levelised one:

\begin{enumerate}
  \item the \textbf{capital recovery factor} is the level annual payment per pound of capex needed to recover it at the target return over the horizon (the standard annuity factor, reducing to one-over-the-horizon at a zero rate);
  \item the \textbf{annual revenue requirement} is the capex times that factor, plus the annual operating cost: the cash the priced structure must collect each year;
  \item the \textbf{fixed price per kWh} is that requirement divided by the annual delivered-energy volume.
\end{enumerate}

The target return itself is a named, versioned assumption. The delivered-energy volume is a dispatch output the current schema does not yet carry, so in this version the price is returned without a value and a warning explains that it awaits dispatch. When the price \emph{is} computed it meets the target by construction, but the model only asserts ``hits target return'' when the target assumption is confirmed; on a placeholder target it computes the price but leaves the ``hits target'' flag unset, because a placeholder cannot certify a return.

\section{The contracted-versus-merchant split}

Lenders draw a hard line between income they can rely on and income they merely hope for. The model surfaces that line directly. The same firmness labels that decided which streams counted toward the priced floor also classify the income for the underwriting pack:

\begin{table}[H]
\centering
\caption{The firm-versus-conditional split of income streams. The contracted floor is what debt is sized against; conditional and gated streams are genuine upside but never underwritten.}
\label{tab:split}
\small\sffamily
\begin{tabularx}{\textwidth}{@{}l l L@{}}
\toprule
\textbf{Firmness} & \textbf{Bucket} & \textbf{Treatment} \\
\midrule
Firm & Contracted floor & counted in priced revenue; debt sized against it \\
Contracted & Contracted floor & counted in priced revenue; debt sized against it \\
Conditional & Merchant upside & held out of the floor; shown as upside only \\
Gated upside (wholesale, balancing, ancillary) & Merchant upside & kept in a separate list; never marked firm; never priced \\
Excluded & Not counted & explicitly priced out \\
\bottomrule
\end{tabularx}
\end{table}

The contracted floor is the bankable case. The merchant upside sits alongside it in the pack as clearly-labelled potential, so a reader sees both the income that supports the loan and the income that might improve the equity return, without the two being blended into one misleadingly rosy figure.

\section{Counterparty credit to rate and tenor (M9)}

The last input is the customer's credit standing. A stronger counterparty is a safer place to lend, so it earns a cheaper, longer loan; a weaker one earns a dearer, shorter loan. This module turns a coarse credit band into the debt rate and tenor the financial model uses, and it does so by a strict, versioned lookup.

The input is a credit \emph{band}, one of \emph{strong}, \emph{adequate} or \emph{weak}, that arrives on the site's counterparty record. (Where that band comes from -- a human assessment now, a credit-data feed later -- is upstream of this module and does not change it; the band is the boundary.) The mapping reads a per-band table from the library that gives each band a rate and a tenor, and emits the matching pair.

\begin{figure}[H]
\centering
\resizebox{\textwidth}{!}{%
\begin{tikzpicture}[node distance=6mm and 13mm]
  \node[io] (band) {credit band\\\scriptsize strong / adequate / weak\\\scriptsize (or unassessed)};
  \node[decide, right=of band] (assessed) {assessed?};
  \node[proc, right=of assessed, text width=30mm] (lookup) {table lookup\\\scriptsize versioned per-band\\\scriptsize rate + tenor};
  \node[decide, right=of lookup, text width=24mm] (mono) {monotonic?};
  \node[io, right=of mono, text width=28mm] (rt) {rate \% + tenor (yr)};
  \node[term, below=8mm of assessed, fill=nodeIO, draw=nodeIOBd] (ask) {data ask:\\get a credit\\assessment};
  \node[term, below=8mm of mono, fill=nodeIO, draw=nodeIOBd] (refuse) {refuse to map:\\policy error};

  \draw[flow] (band) -- (assessed);
  \draw[flow] (assessed) -- node[flowlabel,above]{yes} (lookup);
  \draw[flow] (assessed.south) to[out=-90,in=90] node[flowlabel,right]{no} (ask.north);
  \draw[flow] (lookup) -- (mono);
  \draw[flow] (mono) -- node[flowlabel,above]{yes} (rt);
  \draw[flow] (mono.south) to[out=-90,in=90] node[flowlabel,right]{no} (refuse.north);
\end{tikzpicture}}
\caption{The deterministic credit mapping. An unassessed counterparty becomes a data ask, never the worst band by default; a non-monotonic policy table is refused outright rather than allowed to produce a flattering loan.}
\label{fig:credit}
\end{figure}

Three rules make this safe:

\begin{itemize}
  \item \textbf{No band is ever invented.} An \emph{unknown} or \emph{not-assessed} counterparty is not quietly defaulted to the worst band (that would fabricate an assessment) nor to a kind one. It produces no rate or tenor and becomes a ranked data ask: ``get a credit assessment''.
  \item \textbf{Monotonicity is enforced.} A worse credit band must never buy a cheaper or longer loan. The model checks that, as the band worsens from strong to weak, the rate never falls and the tenor never lengthens. A policy table that violates this is \emph{refused} -- the mapping returns nothing rather than emit, in the brief's words, ``a glowing return at a rate the covenant does not justify''.
  \item \textbf{The numbers are versioned credit policy.} The rate/tenor table lives in the library, not in code, and is placeholder policy until confirmed, so any credit-driven case is marked not bankable until those values are signed off. The mapping is a pure function of the band and the library: rounded at parse time, it reproduces identical debt terms and an identical provenance hash for identical inputs.
\end{itemize}

The mapping also notes, but does not silently fix, a debt tenor that runs longer than the counterparty's expected contract term, because a loan should not outlast the income that services it. That auto-capping is a financing-policy decision deferred until it can be sourced from the library.

\section{Assembling the case, and why it is never bankable today}

Finally the model packs everything into a \code{FinancialCase}: the two cashflow series, the DSCR summary, the two IRRs, payback, the debt assumptions actually used, the price recommendation, the covenant verdict and the deduplicated warnings. It stamps the run with a provenance context (a hash of the timestamp-free inputs plus the assumption, schema and code versions) so the same request reproduces the same case, and the case identifier is derived from that hash rather than from the clock.

The case is marked \textbf{bankable only if every one of four conditions holds}: the cashflows are bankable (a bankable-candidate result built on no unconfirmed assumption), the DSCR is actually available, the covenant passes, and the covenant assumption is itself confirmed. On the current placeholder assumption sets this is \emph{always false}, and the model says so explicitly: the early releases carry no gearing at all (so the DSCR is unavailable), and \code{gb\_v1\_1}, whose placeholder gearing lets the DSCR compute and the covenant gate genuinely fire, is built entirely of labelled placeholders, so the first condition can never hold on it. That is the platform's safety property at work, not a defect: placeholder numbers must never masquerade as bankable ones.

\begin{rebuild}
To re-implement the financial model and credit mapping from scratch:
\begin{enumerate}
  \item \textbf{Gate.} Read the chosen power from the result. If it is missing or non-positive, raise an ``insufficient inputs'' error; the caller turns that into a ranked data ask. Never invent a system.
  \item \textbf{Revenue.} Sum the per-stream annual revenue for streams whose firmness is firm or contracted only. Hold conditional streams and the separate gated-upside list out of the floor. Warn if the floor is empty (cost-only case).
  \item \textbf{Capex.} power (MW) $\times$ installed-cost rate (\pounds/MW, from the library) $\times$ scenario cost multiplier. Reject negatives. Warn (never block, never alter) if the unit cost is outside the library's benchmark band.
  \item \textbf{Cost lines.} Build each line in Table~\ref{tab:costlines} from its single named assumption (levies from the site tariff). A line with a missing driver is \emph{unavailable} with a warning, never zero, never guessed.
  \item \textbf{P50/P90.} Per operating year: P50 revenue $=$ basis revenue $\times$ price multiplier $\times$ $(1-\text{fade})^{\,year-1}$; P90 revenue $=$ P50 $\times$ $(1-\text{haircut})$. Costs identical in both. EBITDA $=$ revenue $-$ dispatch cost $-$ operating cost. Prefix year 0 $= -$capex. Invariant: P90 $\le$ P50 always.
  \item \textbf{Debt + DSCR.} Need gearing, rate, tenor. Principal $=$ capex $\times$ gearing. Debt service by the recorded sculpting basis: level-payment annuity (default), or a P90-DSCR sculpt -- solve the flat target $x$ with PV(P90 cashflow $\div$ $x$, debt rate, tenor) $=$ principal and set service$_y$ $=$ P90 cashflow$_y$ $\div$ $x$; fall back to the annuity with a warning when a P90 year within the tenor is non-positive or the horizon is shorter than the tenor, and record the basis actually applied. DSCR per year $=$ P90 cashflow $\div$ debt service (skip zero-service years). Minimum DSCR is the headline. Compare to the covenant for a pass/fail. If any of gearing/rate/tenor is missing, report DSCR \emph{unavailable} with a naming warning; never guess.
  \item \textbf{Returns.} Payback $=$ interpolated zero-crossing of cumulative P50 cashflow (``does not pay back'' if it never crosses). Project IRR by bisection \emph{only} for a single-sign-change series. Equity IRR on P50 net of the exact debt service from the DSCR step, same reliability gate.
  \item \textbf{Price.} For fixed \pounds/kWh: revenue requirement $=$ capex $\times$ capital-recovery-factor(target return, horizon) $+$ annual opex; price $=$ requirement $\div$ annual delivered energy. Compute only when target return, capex, opex, horizon and delivered energy are all present. Only assert ``hits target'' on a confirmed target. Savings-share and capacity-fee return shape but no value.
  \item \textbf{Credit map.} Map strong/adequate/weak to (rate, tenor) from the versioned per-band table. Unassessed $\rightarrow$ data ask, never a default band. Enforce monotonicity (worse band: rate non-decreasing, tenor non-increasing); refuse a non-monotonic table. Round at parse time for determinism.
  \item \textbf{Assemble + stamp.} Pack the \code{FinancialCase}; stamp a timestamp-free provenance hash with the assumption / schema / code versions. Mark bankable only if cashflows bankable AND DSCR available AND covenant passes AND covenant assumption confirmed -- otherwise not bankable, with a warning.
  \item \textbf{Test.} The case validates; both series present and P90 $\le$ P50; DSCR unavailable when debt inputs are missing; placeholder assumptions raise warnings; outputs scale when a library value changes (no hidden constants); identical inputs give an identical case (determinism); negative capex rejected; the credit map refuses a non-monotonic table and never defaults an unassessed band.
\end{enumerate}
\end{rebuild}

\begin{keyidea}[The one thing to remember]
This module is pure arithmetic with two firm reflexes. It \emph{underwrites to the downside}: debt is sized and covered on the P90 case while returns are reported on P50. And it \emph{degrades honestly}: a missing input becomes an explicit unavailable field and a warning, an unscored counterparty becomes a data ask, a policy that would flatter a weak borrower is refused, and a run built on placeholders is never marked bankable. No AI authors any number here, and no number is ever invented to fill a gap.
\end{keyidea}
